\documentclass[11pt]{article}
\usepackage[margin=0.65in]{geometry}
\usepackage{lmodern}
\usepackage{microtype}
\usepackage{enumitem}
\usepackage{tabularx}
\usepackage[table]{xcolor}
\usepackage{titlesec}
\usepackage[hidelinks]{hyperref}
\setlength{\parindent}{0pt}
\setlength{\parskip}{0.27em}
\setlist[itemize]{leftmargin=1.3em,itemsep=0.07em,topsep=0.1em}
\titleformat{\section}{\large\bfseries}{}{0pt}{}
\titlespacing*{\section}{0pt}{0.6em}{0.18em}
\pagestyle{plain}
\newcommand{\blankline}{\par\vspace{0.1em}\noindent\rule{\linewidth}{0.4pt}\par\vspace{0.34em}}
\newcommand{\smallblank}{\rule{0.43\linewidth}{0.4pt}}
\begin{document}
\begin{center}
{\Large\bfseries U1L15SP --- Short Paper 3}\par
{\LARGE\bfseries What Makes a Man?}\par
{\large Week 15: Verbal, Real, and Mixed Disputes}
\end{center}
\noindent\textbf{Name:} \smallblank\hfill
\textbf{Date:} \rule{1.35in}{0.4pt}

\begingroup\small
\vspace{0.2em}
\fcolorbox{black}{gray!8}{\begin{minipage}{0.94\linewidth}
\textbf{The question}

What does it mean to be a man in \textit{Macbeth}? By the competing meanings used by Macbeth and Lady Macbeth, which character better deserves the name?

You may judge Macbeth, Lady Macbeth, both in different respects, or neither. You must identify the meanings you are applying and test them against the play. The goal is a defensible, qualified judgment---not a teacher-approved definition of masculinity.

This is a question about qualities the characters call manly, not biological sex. Do not assume that violence is the correct definition.
\end{minipage}}

\section*{Text support}
The accompanying \textit{U1L15SP Text Supplement} prints Lady Macbeth's invocation in Act 1, Scene 5 and the manhood dispute in Act 1, Scene 7. These are lines worth focusing on, not limits on your evidence. You may use any accurate line, action, choice, or consequence from the play; remembered events may be paraphrased accurately.

\section*{The assignment}
Write approximately \textbf{two handwritten pages} on separate lined paper. Give it the title \textit{What Makes a Man?} Your paper should:
\begin{itemize}
\item give your own answer to what \emph{man} or \emph{manhood} means in this dispute;
\item explain at least two competing meanings used or implied by Macbeth and Lady Macbeth;
\item decide whether their disagreement is verbal, real, or mixed, and explain what remains disputed after the word is clarified;
\item use at least two brief quotations from the supplement and accurate evidence from later choices or consequences;
\item make a comparative judgment and include a meaningful qualification or counterpressure.
\end{itemize}
Do not force the paper into five paragraphs. Organize it around the definitions, the evidence that tests them, and your final comparison.

\section*{The closed-world rule}
You may use only the supplied text supplement, the assigned play or course-provided primary text, course handouts, and notes you made during class.

\textbf{Do not use AI, the internet, online summaries, outside criticism, a tutor, a parent, a friend, or anyone else's planning, revising, correcting, sentences, or ideas.}

At the end of your paper, write and sign: \emph{I wrote this paper without outside help.}

\section*{What counts --- 12 points}
\rowcolors{2}{gray!8}{white}
\begin{tabularx}{\linewidth}{>{\bfseries}p{0.27\linewidth}Xr}
Claim and terms & Gives a clear judgment and defines the contested term precisely enough to test. & 3\\
Dispute analysis & Explains competing meanings and what real disagreement remains. & 3\\
Textual evidence & Uses accurate quotations and broader-play evidence to test both characters. & 3\\
Explanation & Connects evidence to definitions and comparative judgment. & 2\\
Completion & Is legible, handwritten, and approximately the assigned length. & 1\\
\end{tabularx}
\endgroup

\newpage
\begin{center}
{\Large\bfseries Planning Scaffold}\par
{\large U1L15SP --- Define, test, compare}
\end{center}
\textbf{This scaffold is intentionally shorter than the earlier papers.} Use only the prompts that help you. Notes and fragments are enough; the paper goes on separate lined paper.

\section*{1. Map the contested term}
In a few words, state one meaning of \emph{man} or \emph{manhood} that Lady Macbeth uses or implies, and one that Macbeth uses or implies. Then write what they would still disagree about even if both meanings were made clear.
\blankline
\blankline
\blankline

Classify the dispute as verbal, real, or mixed. What makes that classification fit?
\blankline
\blankline

\section*{2. Test each meaning against the play}
Choose one exact quotation from each supplied passage. For each quotation, note whose meaning it reveals and what it does---and does not---show.
\blankline
\blankline
\blankline

Name one later action or consequence for Macbeth and one for Lady Macbeth that puts a definition under pressure. Paraphrase accurately; do not invent quotations.
\blankline
\blankline

\section*{3. Make a fair comparison}
Write your present judgment in one sentence. Then write the strongest evidence or objection that complicates it.
\blankline
\blankline

What qualification makes your judgment fairer: \emph{in this sense}, \emph{at first}, \emph{although}, \emph{neither}, \emph{both}, or another phrase?
\blankline

\section*{4. Build the paper}
Choose your own order. A useful route is: define the dispute; test Lady Macbeth; test Macbeth; compare; qualify. Change that order if another structure makes the reasoning clearer.

\vfill
\textbf{Before submitting:} Underline your main judgment once. Circle your two exact quotations. Sign the authorship statement.
\end{document}
